\section{Parameters and an exact finite certificate}
\label{sec:parameters}

The preceding counting bounds leave a finite numerical choice: the
partition exponent must be large enough for the archimedean cuts, but
small enough that the number of boxes does not consume the saving in
each box.  We now make that choice.  We also control the simultaneous
increase of the two coordinate weights caused by an unusually short
lattice vector.  The finite calculation below concerns these numerical
inequalities; the geometric estimates supplying the counting thresholds
were proved in the preceding sections.

Throughout this section, $4\le d\le8$, $k=d-2$, and
\[
 K_0=3000,\qquad G_0=100000,\qquad \tau=\frac1{50000}.
\]
Recall the homogeneous polynomial $R(A,B)$ from
\eqref{eq:hilbert-polynomial}, and its associated one-variable
polynomials
\[
 H(t)=d!R(1,t),\qquad
 H_a(t)=(d-1)!(\partial_A R)(1,t),\qquad
 H_b(t)=(d-1)!(\partial_B R)(1,t).
\]
The argument $t$ here is a scalar ratio of bidegrees; it is distinct
from the normalized chart coordinates used earlier.

For an integer $z\ge K_0$, set
\begin{equation}
 \eta_- =\frac z{K_0},\quad
 \eta_+ =\frac{z+1}{K_0},\quad
 b_-=d-k\eta_+,\quad b_+=d-k\eta_-.
 \label{eq:parameter-interval}
\end{equation}
We use these intervals until the first index $z_d$ for which
\begin{equation}
 \eta_-(1-k\eta_-/d)<\frac{2495}{10000}.
 \label{eq:parameter-stop}
\end{equation}
An interval is in mode $\mathsf Q$ when $d=4$ and $z<4000$.
This mode records the use of the stronger quartic surface estimate.

We collect the thresholds that enter the affine counting theorem.
For $U,V,D>0$ and $h\ge2$, write
\[
 \Psi_h(U,V;D)
 =\frac{UV}{\bigl((kU+V)^{h-1}D\bigr)^{1/h}},
 \qquad g=(d-1)(k-1)+\mathbf1_{d=4}.
\]
Keep the mode fixed when evaluating the following functions:
\begin{equation}
\begin{aligned}
 \delta_d(U,V)&=\Psi_d(U,V;U+V),\\
 \delta_h(U,V)&=\Psi_h(U,V;\min(U,V))
       &&(2\le h\le d-1,\ \text{outside mode }\mathsf Q),\\
 \delta_3(U,V)&=\Psi_3(U,V;V)
       &&(\text{mode }\mathsf Q),\\
 \delta_2(U,V)&=\max\left\{
       \Psi_2(U,V;\min(3U,2V,U+V)),\frac U{\sqrt3}\right\}
       &&(\text{mode }\mathsf Q),\\
 \delta_1(U,V)&=\max\left\{\frac U2,\min\left(U,\frac Vg\right)\right\}.
\end{aligned}
\label{eq:parameter-thresholds}
\end{equation}
Finally put $E_h(U,V)=\max_{h\le s\le d}\delta_s(U,V)$.

\begin{proposition}[Parameter choice]
\label{prop:parameter-certificate}
For each $d\in\{4,5,6,7,8\}$, the stopping index in
\eqref{eq:parameter-stop} is the one listed in
Table~\ref{tab:parameter-certificate}.  For every integer
$K_0\le z<z_d$, there are positive rational numbers $t_*,w$ such that,
with $m=(d-1)w$,
\begin{equation}
\begin{gathered}
 0<m<1,\qquad \min(\eta_-,b_-)>m,\\
 w^{d-1}H(t_*)>
 \left(\frac{(d+1)(\eta_++b_+t_*)}{d^2}\right)^d,\\
 w^{d-2}\min\{H_a(t_*),H_b(t_*)\}>
 \left(\frac{\eta_++b_+t_*}{d-1}\right)^{d-1}.
\end{gathered}
\label{eq:parameter-cut-inequalities}
\end{equation}
In particular, the cut conditions \eqref{eq:cut-conditions} hold.
For
\[
 u=\frac{\eta_+-m}{d},\qquad
 v=\frac{b_+-m}{d},
\]
these choices also satisfy
\begin{equation}
 \frac3{20}\le\frac vu\le5,\qquad
 m+\sum_{h=1}^dE_h(u,v)<\frac{999}{1000},
 \qquad
 m+2u<\frac{999}{1000}\quad(\text{mode }\mathsf Q).
 \label{eq:parameter-saving}
\end{equation}
Moreover, for either fixed mode, every $U,V>0$ with
$3/20\le V/U\le5$ satisfies
\begin{equation}
 \delta_h(U+s,V+s)\le\delta_h(U,V)+s
 \quad(s\ge0,\ 1\le h\le d).
 \label{eq:parameter-shift}
\end{equation}
\end{proposition}

\begin{table}[ht]
\centering
\begin{tabular}{c|r|r|c|c}
 $d$ & $z_d$ & Intervals & $z_d/K_0$ & Certified gap before $d\tau$\\
 \hline
 4 & 5124 & 2124 & $427/250$ & $191/100000$\\
 5 & 4084 & 1084 & $1021/750$ & $293/50000$\\
 6 & 3552 & 552 & $148/125$ & $321/2500$\\
 7 & 3226 & 226 & $1613/1500$ & $23177/100000$\\
 8 & 3003 & 3 & $1001/1000$ & $31683/100000$
\end{tabular}
\caption{The intervals have indices $3000\le z<z_d$.
The last column is the minimum of $1-m-\sum_h\overline E_h$,
where each radical threshold is rounded strictly upwards as described
in the proof.  It is therefore a lower bound for the true gap.}
\label{tab:parameter-certificate}
\end{table}

\begin{proof}
We first give the rational prescription for $t_*,w$, together with
an exact check of the finite inequalities.  We then prove the
continuous assertion \eqref{eq:parameter-shift}.

\paragraph{Polynomial coefficients and the choice of weights.}
If
\[
 c_j=\binom dj\binom{d-2}{j-1}k^{j-1}\qquad(1\le j\le d-1),
\]
then direct differentiation of $R$ gives
\begin{equation}
 H(t)=\sum_{j=1}^{d-1}c_jt^j,\quad
 H_a(t)=\sum_{j=1}^{d-1}\frac{d-j}{d}c_jt^j,\quad
 H_b(t)=\sum_{j=1}^{d-1}\frac jdc_jt^{j-1}.
 \label{eq:parameter-polynomials}
\end{equation}
The binomial identities
$(d-j)\binom dj/d=\binom{d-1}j$ and
$j\binom dj/d=\binom{d-1}{j-1}$ explain the three sums used in the
certificate below.

Choose $t_*$ among $i/50$, $1\le i<500$, to minimize
\[
 \frac{(\eta_++b_+t)^d}{H(t)},
\]
breaking a tie by taking the first value.  For $0\le x<1$ and an
integer $e\ge1$, let $\operatorname{up}_e(x)$ be the least grid
point $r/G_0$, with $r\in\{1,\ldots,G_0\}$, whose $e$-th power is
strictly larger than $x$.
Set
\begin{equation}
 w=\operatorname{up}_{d-1}\left(
 \frac{1}{H(t_*)}
 \left(\frac{(d+1)(\eta_++b_+t_*)}{d^2}\right)^d
 \right).
 \label{eq:parameter-prescription}
\end{equation}
The certificate checks that the argument lies in $[0,1)$.
The first strict cut inequality then follows from the definition of
$\operatorname{up}_{d-1}$; the second one and all the other finite
conditions are checked directly.

There is no floating-point root evaluation.  The function
\texttt{upper} in the complete Python certificate of
Appendix~\ref{app:certificate} maintains
\[
 (l/G_0)^e\le x<(r/G_0)^e
\]
and terminates when $r=l+1$.  Thus the output is a strict upper
bound, including when the exact root is already a grid point.
Replacing each radical in \eqref{eq:parameter-thresholds} by such an
upper bound can only increase each $E_h$.  The list \texttt{P}
in that certificate is ordered from dimension $d$ down to dimension $1$; its
successive prefix maxima therefore give exactly these conservative
bounds for $E_d,E_{d-1},\ldots,E_1$.



All comparisons in this certificate use integers or rational numbers.
It verifies every interval and the entries of
Table~\ref{tab:parameter-certificate}, rather than an interpolation
between sampled intervals.  The choice of the two endpoint weights
in \eqref{eq:parameter-interval} makes each successful check valid
throughout its full interval.

\paragraph{The common-shift estimate.}
For $a,b\ge0$ with $a+b>0$, define
\[
 q_{a,b}(U,V)=\frac{UV}{aU+bV},\qquad
 F_{h,a,b}(U,V)=\frac{UV}
 {\bigl((kU+V)^{h-1}(aU+bV)\bigr)^{1/h}}.
\]
The Hessian of $q_{a,b}$ is
\[
 \frac{-2ab}{(aU+bV)^3}
 \begin{pmatrix}V^2&-UV\\-UV&U^2\end{pmatrix},
\]
so $q_{a,b}$ is concave.  It is also increasing in each variable;
when one coefficient is zero it is a positive multiple of a
coordinate.  The identity
\[
 F_{h,a,b}=q_{k,1}^{(h-1)/h}q_{a,b}^{1/h}
\]
and concavity and monotonicity of the weighted geometric mean show
that each $F_{h,a,b}$ is concave and increasing.  It is homogeneous
of degree one.

Write $f(t)=F_{h,a,b}(1,t)$.  Homogeneity shows that the derivative
in the direction $(1,1)$, at a point with ratio $t=V/U$, is
\[
 D(t)=f(t)+(1-t)f'(t).
\]
Since $D'(t)=(1-t)f''(t)$, this derivative is nonincreasing before
$t=1$ and nondecreasing after $t=1$.  Its maximum on $[3/20,5]$
is attained at an endpoint.  Direct differentiation gives
\[
 D(t)=\frac{N_{h,a,b}(t)}
 {\bigl((k+t)^{h-1}(a+bt)\bigr)^{1/h}},
\]
where
\[
 N_{h,a,b}(t)=t\left(1+\frac1t
 -\frac{h-1}{h}\frac{k+1}{k+t}
 -\frac1h\frac{a+b}{a+bt}\right).
\]
The function \texttt{check\_derivatives} verifies, for both endpoints,
all $2\le k\le6$, all $2\le h\le k+2$, and
$(a,b)=(1,0),(0,1),(1,1)$, that
\[
 N_{h,a,b}(t)\ge0,\qquad
 N_{h,a,b}(t)^h<(k+t)^{h-1}(a+bt).
\]
Thus $D(t)\le1$ on the full ratio interval for these branches.
The shifted ratio
\[
 \frac{V+s}{U+s}
 =\frac U{U+s}\frac VU+\frac s{U+s}
\]
remains between $V/U$ and $1$.  Integrating the derivative bound
along this path gives the claimed shift estimate for every branch.

Multiplying the denominator form by a constant $c\ge1$ multiplies
the branch by $c^{-1/h}\le1$, so the forms $3U$ and $2V$ are
covered as well.  A finite minimum or maximum of functions satisfying
the shift estimate satisfies the same estimate, simply by taking
that extremum on both sides.  This handles the min/max expressions
in \eqref{eq:parameter-thresholds}; no concavity of their maximum
is needed.  Finally, the functions $U/\sqrt3$, $U/2$, $U$, and
$V/g$ have directional slopes at most one.  This also proves the
curve case and completes \eqref{eq:parameter-shift}.
\end{proof}

\subsection{Coverage and the remaining exponent margins}

The transition function $\phi_d(\eta)=\eta(1-k\eta/d)$ is
nonincreasing for $\eta\ge1$, since
\[
 \phi_d'(\eta)=1-\frac{2k\eta}{d}
 \le\frac{4-d}{d}\le0.
\]
Consequently all intervals before the stopping index are covered by
Proposition~\ref{prop:parameter-certificate}, and every subsequent
$\eta\le d/k$ satisfies the larger-$\eta$ condition
$\phi_d(\eta)<2495/10000$.  The last closed interval ends at
$z_d/K_0$, so the two ranges meet without a gap.  In the quartic case,
the last $\mathsf Q$ interval ends at $4/3$ and the next interval
uses the ordinary thresholds; either interval covers their common
endpoint.

The lattice ranges of Proposition~\ref{prop:lattice-boxes} have
$U=u+j+\tau$, $V=v+j+\tau$ and box factor $X^{m-dj}$.
By \eqref{eq:parameter-shift}, their combined exponent is at most
\begin{equation}
 m-dj+\sum_{h=1}^dE_h(U,V)
 \le m+\sum_{h=1}^dE_h(u,v)+d\tau
 <\frac{24979}{25000}<1.
 \label{eq:parameter-final-box-exponent}
\end{equation}
The estimate uses only the advertised $999/1000$ bound and $d\le8$;
Table~\ref{tab:parameter-certificate} gives stronger bounds.  For
mode $\mathsf Q$, the removed base fibers have combined exponent
\begin{equation}
 m-4j+2U=m+2u-2j+2\tau
 <\frac{3122}{3125}<1.
 \label{eq:parameter-fiber-exponent}
\end{equation}
In particular, both inequalities leave positive room for the
arbitrarily small loss in the affine counting theorem.

The direct larger-$\eta$ argument uses a fixed
finite set of positive rational coordinate weights even at the
terminal point $b=d-k\eta=0$.  The available product gap is
\begin{equation}
 \left(\frac{4999}{10000}\right)^2-\frac{2495}{10000}
 =\frac{40001}{100000000}>0.
 \label{eq:parameter-large-eta-gap}
\end{equation}
An explicit choice on $1\le\eta\le d/k$ is, with $e=1/10000$,
\[
 b'=e(\lfloor b/e\rfloor+1),\qquad
 \eta'=e(\lfloor\eta/e\rfloor+1).
\]
These are strict positive upper weights, drawn from a finite set.
Since $b+\eta\le3$ and $d\ge4$, whenever
$b\eta/d\le2495/10000$ they satisfy
\[
 \frac{b'\eta'}d
 \le\frac{2495}{10000}+\frac{3e+e^2}{4}
 <\left(\frac{4999}{10000}\right)^2.
\]
Also include the pair $(b',\eta')=(e,d/k+e)$, which covers
$d/k\le\eta<d/k+e$.  This includes every fixed-coefficient
$O_f(1/\log X)$ overshoot once $X$ is sufficiently large.  Thus the
surface threshold in the direct argument is strictly below
$4999/10000$, including the endpoint range.  Together with the curve
threshold $1/2$, the exponent from the two cumulative maxima is
strictly below $9999/10000$.  These fixed margins are the numerical
input for the final tail summation.
